% !TEX root = ../main.tex
\chapter{Kesimpulan dan Saran}

\section{Kesimpulan}

Penelitian ini menjawab tiga rumusan masalah yang diajukan pada Bab~I.

Pertama, pemanfaatan \textit{Generative AI} sebagai anotator terbukti mampu menstandarkan label aktivitas pada \textit{event log} pemeliharaan melalui abstraksi semantik berbasis taksonomi dua belas kelas. \textit{Large Language Model} \texttt{gemini-3.1-flash-lite} memetakan deskripsi mentah yang ditulis bebas oleh teknisi menjadi kategori tindakan yang seragam, dan seluruh 308 \textit{event} berhasil dianotasi tanpa bergantung pada aturan berbasis kata kunci.

Kedua, perbandingan \textit{process discovery} pada \textit{event log} \textit{as-is} dan \textit{enriched} menghasilkan dua bentuk model proses yang sangat berbeda. Model \textit{as-is} yang dibangun dari 112 label mentah menghasilkan struktur padat berisi 81 \textit{place} dan 119 \textit{transition}, sedangkan model \textit{enriched} yang dibangun dari 12 kategori tindakan menghasilkan struktur ringkas berisi 31 \textit{place} dan 43 \textit{transition}. Abstraksi semantik mengubah model yang menyerupai spaghetti menjadi alur pemeliharaan yang dapat dibaca dengan jelas.

Ketiga, penyeragaman label aktivitas berbasis LLM menghasilkan model yang lebih sederhana dan mudah dipahami dibandingkan metode konvensional tanpa anotasi. Kesederhanaan ini paling kuat terlihat pada pengurangan aktivitas unik sebesar 89,29\% serta penyusutan \textit{place} dan \textit{transition} masing-masing sebesar 61,73\% dan 63,87\%. Evaluasi \textit{conformance} memperkuat temuan ini, yaitu \textit{generalization} meningkat tajam dari 0,0807 menjadi 0,7079 sementara \textit{fitness} tetap tinggi pada 0,9570. Penurunan \textit{precision} dari 0,8716 menjadi 0,5063 mencerminkan berkurangnya \textit{overfitting} model \textit{as-is}, bukan kemunduran kualitas, sebab tujuan abstraksi semantik adalah menghasilkan model yang dapat digeneralisasi dan bermakna secara operasional, bukan model yang menghafal setiap \textit{trace}.

\section{Saran}

Beberapa arah dapat memperkuat dan memperluas penelitian ini. Cakupan data perlu diperluas ke unit pembangkit lain dan jenis pekerjaan yang berbeda, mengingat data penelitian ini berasal dari satu unit PLTD sehingga nilai temuannya terbatas pada lingkup tersebut, sementara \textit{pipeline}-nya berpotensi digeneralisasi. Penelitian lanjutan juga dapat membandingkan beberapa model LLM serta menguji sensitivitas hasil terhadap perubahan taksonomi dan parameter \textit{temperature}. Terakhir, integrasi metrik keterbacaan yang lebih representatif daripada \textit{simplicity} PM4Py akan membantu mengukur kemudahan interpretasi model secara lebih adil.
